\chapter{Boolean Analysis --- Normalized Sums}
%%% generated by the blueprint pipeline from TCSlib.BooleanAnalysis.GaussianSpace.NormalizedSums
\section{Overview}
The multilinear polynomial for a normalized sum has degree one, equal coordinate influences, and unit variance in positive dimension. Moment-matched independent inputs admit quantitative Gaussian comparison against Lipschitz tests.

\section{Declarations}

\begin{definition}[Normalized-sum multilinear polynomial]
\lean{BooleanAnalysis.GaussianSpace.normalizedSumPolynomial}
\leanok
\uses{BooleanAnalysis.ThresholdFunctions.MultilinearPolynomial}
For $n\ge0$, define a multilinear coefficient family $P_n$ on subsets of $\{1,\ldots,n\}$ by $P_n(S)=1/\sqrt n$ when $|S|=1$ and $P_n(S)=0$ otherwise; division by zero is zero. Its evaluation is $n^{-1/2}\sum_i z_i$.
\end{definition}

\begin{lemma}[Degree of the normalized-sum polynomial]
\lean{BooleanAnalysis.GaussianSpace.normalizedSumPolynomial_hasDegreeAtMost}
\leanok
\uses{BooleanAnalysis.GaussianSpace.normalizedSumPolynomial,BooleanAnalysis.ThresholdFunctions.MultilinearPolynomial.HasDegreeAtMost}
\difficulty{1}
For every $n\ge0$, the multilinear polynomial with singleton coefficients $1/\sqrt n$ and all other coefficients zero has degree at most one.
\end{lemma}

\begin{lemma}[Influence of a normalized-sum coordinate]
\lean{BooleanAnalysis.GaussianSpace.influence_normalizedSumPolynomial}
\leanok
\uses{BooleanAnalysis.GaussianSpace.normalizedSumPolynomial,BooleanAnalysis.ThresholdFunctions.MultilinearPolynomial.eval,BooleanAnalysis.ThresholdFunctions.MultilinearPolynomial.fourierCoeff_eval,BooleanAnalysis.influence,BooleanAnalysis.influence_eq_sum_fourier}
\difficulty{5}
For $n\ge0$ and a coordinate $i\in\{1,\ldots,n\}$, let $f(x)=n^{-1/2}\sum_j(-1)^{x_j}$ on the uniform Boolean cube. Its influence $\mathbb E(D_i f)^2$, with $D_i$ the coordinate half-difference, equals $1/n$.
\end{lemma}

\begin{lemma}[Variance of the normalized-sum polynomial]
\lean{BooleanAnalysis.GaussianSpace.cubeVariance_normalizedSumPolynomial}
\leanok
\uses{BooleanAnalysis.GaussianSpace.influence_normalizedSumPolynomial,BooleanAnalysis.GaussianSpace.normalizedSumPolynomial,BooleanAnalysis.Hypercontractivity.cubeVariance,BooleanAnalysis.Hypercontractivity.cubeVariance_eq_sum_nonempty_fourier,BooleanAnalysis.ThresholdFunctions.MultilinearPolynomial.eval,BooleanAnalysis.ThresholdFunctions.MultilinearPolynomial.fourierCoeff_eval,BooleanAnalysis.totalInfluence,BooleanAnalysis.totalInfluence_eq_sum_sq_deg}
\difficulty{5}
If $n>0$ and $f(x)=n^{-1/2}\sum_{i=1}^n(-1)^{x_i}$ on the uniform Boolean cube, then $\operatorname{Var}(f)=1$.
\end{lemma}

\begin{lemma}[Evaluation at arbitrary random inputs]
\lean{BooleanAnalysis.GaussianSpace.evalMultilinearRV_normalizedSumPolynomial}
\leanok
\uses{BooleanAnalysis.GaussianSpace.normalizedSumPolynomial,BooleanAnalysis.Hypercontractivity.evalMultilinearRV}
\difficulty{5}
Let $n\ge0$, let $\Omega$ be a measurable space, and let $X_i:\Omega\to\mathbb R$ be arbitrary functions. The multilinear polynomial with singleton coefficients $1/\sqrt n$ and all other coefficients zero evaluates at $(X_i(\omega))_i$ to $\sum_iX_i(\omega)/\sqrt n$ for every $\omega\in\Omega$, with division by zero equal to zero.
\end{lemma}

\begin{lemma}[Gaussian comparison for independent normalized sums]
\lean{BooleanAnalysis.GaussianSpace.normalizedSum_lipschitz_comparison}
\leanok
\uses{BooleanAnalysis.GaussianSpace.HasInvarianceMoments,BooleanAnalysis.GaussianSpace.HasInvarianceMoments.identDistrib,BooleanAnalysis.GaussianSpace.cubeVariance_normalizedSumPolynomial,BooleanAnalysis.GaussianSpace.evalMultilinearRV_normalizedSumPolynomial,BooleanAnalysis.GaussianSpace.hasInvarianceMoments_standardGaussian,BooleanAnalysis.GaussianSpace.influence_normalizedSumPolynomial,BooleanAnalysis.GaussianSpace.invariance_lipschitz,BooleanAnalysis.GaussianSpace.measurePreserving_gaussianProjection,BooleanAnalysis.GaussianSpace.normalizedSumPolynomial,BooleanAnalysis.GaussianSpace.normalizedSumPolynomial_hasDegreeAtMost,BooleanAnalysis.ThresholdFunctions.standardGaussianMeasure}
\difficulty{6}
There is a universal $K>0$ such that, for every $n>0$, probability space $(\Omega,\mu)$, and independent real random variables $X_1,\ldots,X_n$ with finite fourth moments, $\mathbb E X_i=\mathbb E X_i^3=0$, $\mathbb E X_i^2=1$, and $\mathbb E X_i^4\le9$, every $c$-Lipschitz $\psi:\mathbb R\to\mathbb R$ with $c\ge0$ satisfies $|\mathbb E\psi(n^{-1/2}\sum_iX_i)-\mathbb E\psi(Z)|\le Kc\,n^{-1/4}$, where $Z$ is standard Gaussian.
\end{lemma}

\begin{lemma}[Product integration equals cube expectation]
\lean{BooleanAnalysis.GaussianSpace.integral_uniformBool_pi_eq_expect}
\leanok
\uses{BooleanAnalysis.BoolCube,BooleanAnalysis.BooleanFunc,BooleanAnalysis.expect,BooleanAnalysis.moment_eq_expect,BooleanAnalysis.uniformWeight}
\difficulty{4}
For every $n\ge0$ and $f:\{0,1\}^n\to\mathbb R$, integration against the product of $n$ uniform bit measures equals the uniform cube expectation $2^{-n}\sum_xf(x)$.
\end{lemma}

\begin{lemma}[Gaussian comparison for normalized Boolean sums]
\lean{BooleanAnalysis.GaussianSpace.normalizedBoolSum_lipschitz_comparison}
\leanok
\uses{BooleanAnalysis.GaussianSpace.HasInvarianceMoments,BooleanAnalysis.GaussianSpace.HasInvarianceMoments.identDistrib,BooleanAnalysis.GaussianSpace.hasInvarianceMoments_uniformBool,BooleanAnalysis.GaussianSpace.normalizedSum_lipschitz_comparison,BooleanAnalysis.ThresholdFunctions.standardGaussianMeasure,BooleanAnalysis.boolToSign}
\difficulty{5}
There is a universal $K>0$ such that, for $n>0$, independent uniform signs $\epsilon_1,\ldots,\epsilon_n$, a standard Gaussian $Z$, and any $c$-Lipschitz $\psi:\mathbb R\to\mathbb R$ with $c\ge0$, $|\mathbb E\psi(n^{-1/2}\sum_i\epsilon_i)-\mathbb E\psi(Z)|\le Kc\,n^{-1/4}$.
\end{lemma}

\begin{lemma}[Gaussian limit for normalized Boolean sums]
\lean{BooleanAnalysis.GaussianSpace.tendsto_integral_normalizedBoolSum}
\leanok
\uses{BooleanAnalysis.GaussianSpace.normalizedBoolSum_lipschitz_comparison,BooleanAnalysis.ThresholdFunctions.standardGaussianMeasure,BooleanAnalysis.boolToSign}
\difficulty{5}
For every $c\ge0$ and $c$-Lipschitz $\psi:\mathbb R\to\mathbb R$, if $\epsilon_1,\ldots,\epsilon_{m+1}$ are independent uniform signs and $Z$ is standard Gaussian, then $\mathbb E\psi((m+1)^{-1/2}\sum_i\epsilon_i)\to\mathbb E\psi(Z)$ as $m\to\infty$, where the first expectation is integration under the product sign law.
\end{lemma}

\begin{lemma}[Cube-expectation limit for normalized sums]
\lean{BooleanAnalysis.GaussianSpace.tendsto_expect_normalizedBoolSum}
\leanok
\uses{BooleanAnalysis.BoolCube,BooleanAnalysis.GaussianSpace.integral_uniformBool_pi_eq_expect,BooleanAnalysis.GaussianSpace.tendsto_integral_normalizedBoolSum,BooleanAnalysis.ThresholdFunctions.standardGaussianMeasure,BooleanAnalysis.boolToSign,BooleanAnalysis.expect}
\difficulty{1}
For every $c\ge0$ and $c$-Lipschitz $\psi:\mathbb R\to\mathbb R$, $2^{-(m+1)}\sum_{x\in\{0,1\}^{m+1}}\psi((m+1)^{-1/2}\sum_i(-1)^{x_i})\to\mathbb E\psi(Z)$ as $m\to\infty$, where $Z$ is standard Gaussian.
\end{lemma}
